Une lumière de \qty{\v{lam}}{nm} éclaire deux plaques. Le métal 1 a un travail d'extraction de \qty{\v{phi1}}{eV} et le métal 2, de \qty{\v{phi2}}{eV}.

\begin{enumerate}[label=\alph*)]
    \item Déterminez le potentiel d'arrêt pour chaque métal.
    \item On augmente graduellement la longueur d'onde de la lumière. À partir de quelle longueur d'onde le métal 2 cesse-t-il d'émettre des électrons ? Le métal 1 en émet-il encore à cette longueur d'onde ?
    \item Si l'on double l'intensité lumineuse sans changer la longueur d'onde, qu'arrive-t-il aux potentiels d'arrêt et au nombre d'électrons éjectés ?
\end{enumerate}
